\documentclass[10pt,a4paper]{amsart}
\usepackage{amsmath,amssymb,amsthm,graphicx}
\usepackage{verbatim}
\newenvironment{reference}{\comment}{\endcomment}
%\newenvironment{reference}{}{}
\newenvironment{slogan}{\comment}{\endcomment}
\newenvironment{history}{\comment}{\endcomment}

% For commutative diagrams we use Xy-pic
\usepackage[all]{xy}

% We use 2cell for 2-commutative diagrams.
\xyoption{2cell}
% Theorem environments.
%
\theoremstyle{plain}
\newtheorem{theorem}[subsection]{Theorem}
\newtheorem{proposition}[subsection]{Proposition}
\newtheorem{lemma}[subsection]{Lemma}

\theoremstyle{definition}
\newtheorem{definition}[subsection]{Definition}
\newtheorem{example}[subsection]{Example}
\newtheorem{exercise}[subsection]{Exercise}
\newtheorem{situation}[subsection]{Situation}

\theoremstyle{remark}
\newtheorem{remark}[subsection]{Remark}
\newtheorem{remarks}[subsection]{Remarks}

\numberwithin{equation}{subsection}

% Macros
%
\def\lim{\mathop{\mathrm{lim}}\nolimits}
\def\colim{\mathop{\mathrm{colim}}\nolimits}
\def\Spec{\mathop{\mathrm{Spec}}}
\def\Hom{\mathop{\mathrm{Hom}}\nolimits}
\def\Ext{\mathop{\mathrm{Ext}}\nolimits}
\def\SheafHom{\mathop{\mathcal{H}\!\mathit{om}}\nolimits}
\def\SheafExt{\mathop{\mathcal{E}\!\mathit{xt}}\nolimits}
\def\Sch{\mathit{Sch}}
\def\Mor{\mathop{\mathrm{Mor}}\nolimits}
\def\Ob{\mathop{\mathrm{Ob}}\nolimits}
\def\Sh{\mathop{\mathit{Sh}}\nolimits}
\def\NL{\mathop{N\!L}\nolimits}
\def\CH{\mathop{\mathrm{CH}}\nolimits}
\def\proetale{{pro\text{-}\acute{e}tale}}
\def\etale{{\acute{e}tale}}
\def\QCoh{\mathit{QCoh}}
\def\Ker{\mathop{\mathrm{Ker}}}
\def\Im{\mathop{\mathrm{Im}}}
\def\Coker{\mathop{\mathrm{Coker}}}
\def\Coim{\mathop{\mathrm{Coim}}}

% Boxtimes
%
\DeclareMathSymbol{\boxtimes}{\mathbin}{AMSa}{"02}

%
% Macros for moduli stacks/spaces
%
\def\QCohstack{\mathcal{QC}\!\mathit{oh}}
\def\Cohstack{\mathcal{C}\!\mathit{oh}}
\def\Spacesstack{\mathcal{S}\!\mathit{paces}}
\def\Quotfunctor{\mathrm{Quot}}
\def\Hilbfunctor{\mathrm{Hilb}}
\def\Curvesstack{\mathcal{C}\!\mathit{urves}}
\def\Polarizedstack{\mathcal{P}\!\mathit{olarized}}
\def\Complexesstack{\mathcal{C}\!\mathit{omplexes}}
% \Pic is the operator that assigns to X its picard group, usage \Pic(X)
% \Picardstack_{X/B} denotes the Picard stack of X over B
% \Picardfunctor_{X/B} denotes the Picard functor of X over B
\def\Pic{\mathop{\mathrm{Pic}}\nolimits}
\def\Picardstack{\mathcal{P}\!\mathit{ic}}
\def\Picardfunctor{\mathrm{Pic}}
\def\Deformationcategory{\mathcal{D}\!\mathit{ef}}
\usepackage{cleveref}
\crefname{originallemma}{original lemma}{original lemmas}
\Crefname{originallemma}{Original lemma}{Original lemmas}
\makeatletter
\newlabel{algebra-lemma-faithfully-flat-universally-injective}{{05CK}{source}{}{sourcepointer-1}{}}
\newlabel{algebra-lemma-faithfully-flat-universally-injective@cref}{{[originallemma][0][]05CK}{[1][0][]source}{}{}{}}
\makeatother
\crefname{originalreference}{original reference}{original references}
\Crefname{originalreference}{Original reference}{Original references}
\makeatletter
\newlabel{lemma-bourbaki-finite-type-general-base}{{05IJ}{source}{}{sourcepointer-2}{}}
\newlabel{lemma-bourbaki-finite-type-general-base@cref}{{[originalreference][0][]05IJ}{[1][0][]source}{}{}{}}
\makeatother
\crefname{originalreference}{original reference}{original references}
\Crefname{originalreference}{Original reference}{Original references}
\makeatletter
\newlabel{lemma-contains-relative-ass-after-base-change}{{0CW2}{source}{}{sourcepointer-3}{}}
\newlabel{lemma-contains-relative-ass-after-base-change@cref}{{[originalreference][0][]0CW2}{[1][0][]source}{}{}{}}
\makeatother
\crefname{originallemma}{original lemma}{original lemmas}
\Crefname{originallemma}{Original lemma}{Original lemmas}
\makeatletter
\newlabel{spaces-divisors-lemma-ses-weakly-ass}{{0CTZ}{source}{}{sourcepointer-4}{}}
\newlabel{spaces-divisors-lemma-ses-weakly-ass@cref}{{[originallemma][0][]0CTZ}{[1][0][]source}{}{}{}}
\makeatother
\crefname{originallemma}{original lemma}{original lemmas}
\Crefname{originallemma}{Original lemma}{Original lemmas}
\makeatletter
\newlabel{spaces-divisors-lemma-weakly-ass-zero}{{0CU0}{source}{}{sourcepointer-5}{}}
\newlabel{spaces-divisors-lemma-weakly-ass-zero@cref}{{[originallemma][0][]0CU0}{[1][0][]source}{}{}{}}
\makeatother
\title{If the associated points lie in the image of a flat morphism then the canonical map to the pushforward is universally injective}
\author{The Stacks Project Authors}
\date{stacks-project:spaces-flat.tex:local-snapshot-20261012; GNU FDL 1.2 or later}
\begin{document}
\maketitle
\noindent\textit{Editorial note.} All delivered mathematical text is the author's own, copied from the pinned Stacks Project chapter file spaces-flat.tex; no statement or proof step was rewritten, reordered or removed.\par\medskip
\noindent\textit{Editorial note.} The counted proof is the author's complete argument: a reduction to weakly associated points of the kernel and a point by point argument using faithful flatness of the local ring map along a flat morphism with preimages of all associated points\par\medskip
\noindent\textit{Editorial note.} The label lemma-universally-separating names the counted statement, and the proof environment immediately following it in the source is delivered verbatim as its proof.\par\medskip
\noindent\textit{Editorial note.} The counted statement is carried by the author's own lemma environment in the chapter source.\par\medskip
\noindent\textit{Allowed standard prerequisites (editorial).}\begin{itemize}
\item Flatness and weak associated points, and the criterion that a module is zero when it has no weakly associated point
\item The inclusion of the weakly associated points of a subquotient into those of the ambient module, and the containment of the weakly associated points of a flat finite type module in the associated points of the fibres
\item Faithful flatness of the map of local rings induced by a flat morphism at a point of the image, and the resulting injectivity of the base changed stalk
\end{itemize}\medskip
\ifdefined\hypertarget\hypertarget{sourcepointer-1}{}\fi
\noindent\textit{Source reference.} Original lemma 05CK: sources/stacks-chow/tags.txt, source lines 6923--6923. Prerequisite pointer invoked at proof line(s) 1208; it names a standard fact of a sibling chapter that the counted argument uses, and every step of the counted proof built on it is retained verbatim in the delivered body.\par\medskip
\ifdefined\hypertarget\hypertarget{sourcepointer-2}{}\fi
\noindent\textit{Source reference.} Original reference 05IJ: sources/stacks-spaces-flat/pinned-original-source.tex, source lines 733--733. Prerequisite pointer invoked at proof line(s) 1193; it names a standard fact of the same chapter that the counted argument uses, and every step of the counted proof built on it is retained verbatim in the delivered body.\par\medskip
\ifdefined\hypertarget\hypertarget{sourcepointer-3}{}\fi
\noindent\textit{Source reference.} Original reference 0CW2: sources/stacks-spaces-flat/pinned-original-source.tex, source lines 1031--1031. Prerequisite pointer invoked at proof line(s) 1180; it names a standard fact of the same chapter that the counted argument uses, and every step of the counted proof built on it is retained verbatim in the delivered body.\par\medskip
\ifdefined\hypertarget\hypertarget{sourcepointer-4}{}\fi
\noindent\textit{Source reference.} Original lemma 0CTZ: sources/stacks-chow/tags.txt, source lines 16084--16084. Prerequisite pointer invoked at proof line(s) 1191; it names a standard fact of a sibling chapter that the counted argument uses, and every step of the counted proof built on it is retained verbatim in the delivered body.\par\medskip
\ifdefined\hypertarget\hypertarget{sourcepointer-5}{}\fi
\noindent\textit{Source reference.} Original lemma 0CU0: sources/stacks-chow/tags.txt, source lines 16085--16085. Prerequisite pointer invoked at proof line(s) 1188; it names a standard fact of a sibling chapter that the counted argument uses, and every step of the counted proof built on it is retained verbatim in the delivered body.\par\medskip
\begin{lemma}
\label{lemma-universally-separating}
Let $Y$ be an algebraic space over a scheme $S$.
Let $g : X' \to X$ be a flat morphism of algebraic spaces over $Y$
with $X$ locally of finite type over $Y$.
Let $\mathcal{F}$ be a finite type quasi-coherent $\mathcal{O}_X$-module
which is flat over $Y$. If $\text{Ass}_{X/Y}(\mathcal{F}) \subset g(|X'|)$
then the canonical map
$$
\mathcal{F} \longrightarrow g_*g^*\mathcal{F}
$$
is injective, and remains injective after any base change.
\end{lemma}

\begin{proof}
The final assertion means that $\mathcal{F}_Z \to (g_Z)_*g_Z^*\mathcal{F}_Z$
is injective for any morphism $Z \to Y$. Since the assumption on
the relative assassin is preserved by base change
(Lemma \ref{lemma-contains-relative-ass-after-base-change})
it suffices to prove the injectivity of the displayed arrow.

\medskip\noindent
Let $\mathcal{K} = \Ker(\mathcal{F} \to g_*g^*\mathcal{F})$.
Our goal is to prove that $\mathcal{K} = 0$.
In order to do this it suffices to prove that
$\text{WeakAss}_X(\mathcal{K}) = \emptyset$, see
Divisors on Spaces, Lemma \ref{spaces-divisors-lemma-weakly-ass-zero}.
We have
$\text{WeakAss}_X(\mathcal{K}) \subset \text{WeakAss}_X(\mathcal{F})$, see
Divisors on Spaces, Lemma \ref{spaces-divisors-lemma-ses-weakly-ass}.
As $\mathcal{F}$ is flat we see from
Lemma \ref{lemma-bourbaki-finite-type-general-base}
that $\text{WeakAss}_X(\mathcal{F}) \subset \text{Ass}_{X/Y}(\mathcal{F})$.
By assumption any point $x$ of $\text{Ass}_{X/Y}(\mathcal{F})$
is the image of some $x' \in |X'|$. Since $g$ is flat the
local ring map
$\mathcal{O}_{X, \overline{x}} \to \mathcal{O}_{X', \overline{x}'}$
is faithfully flat, hence the map
$$
\mathcal{F}_{\overline{x}}
\longrightarrow
(g^*\mathcal{F})_{\overline{x}'} =
\mathcal{F}_{\overline{x}} \otimes_{\mathcal{O}_{X, \overline{x}}}
\mathcal{O}_{X', \overline{x}'}
$$
is injective (see
Algebra, Lemma \ref{algebra-lemma-faithfully-flat-universally-injective}).
Since the displayed arrow factors through
$\mathcal{F}_{\overline{x}} \to (g_*g^*\mathcal{F})_{\overline{x}}$,
we conclude that
$\mathcal{K}_{\overline{x}} = 0$. Hence $x$ cannot be a weakly associated
point of $\mathcal{K}$ and we win.
\end{proof}

\end{document}
